%% kapa: exam=1 round=30 year=2019 period=1 form=A subject=civil_law subject_ko=민법 seq=22 number=22
%% source: Q-Net 감정평가사 기출문제 (cst00309) artlSeq=5205205 file=extracted/1008/5205205/2019년 제30회 감정평가사 1차시험 문제지/2019년 제30회 감정평가사 1차시험 문제지/2019년도 제30회 감정평가사 1차 1교시 A형.hwp
%% conversion: kordoc 4.24.51 -> tools/gen1.py (auto); review=proofread-sonnet; flags=none
\begin{question}{22}
물권적 청구권에 관한 설명으로 옳은 것을 모두 고른 것은? (다툼이 있으면 판례에 따름)
\begin{framedtext}
ㄱ. 부동산 매매계약이 합의해제되면 매수인에게 이전되었던 소유권은 당연히 매도인에게 복귀되므로 합의해제에 따른 매도인의 원상회복청구권은 소유권에 기인한 물권적 청구권으로서 이는 소멸시효의 대상이 아니다. \par
ㄴ. 임대차목적물 침해자에 대하여 임차인은 점유보호청구권을 행사할 수 있으나, 소유자인 임대인은 점유보호청구권을 행사할 수 없다. \par
ㄷ. 불법한 원인으로 급여를 한 사람은 그 원인행위가 법률상 무효라 하여 상대방에게 부당이득반환청구를 할 수 없는 경우, 급여한 물건의 소유권이 여전히 자기에게 있다고 하여 소유권에 기한 반환청구도 할 수 없다. \par
ㄹ. 물건의 양도 시 소유권에 기한 물권적 청구권을 소유권과 분리하여 이를 소유권을 상실한 전(前)소유자에게 유보하여 행사시킬 수 있다.
\end{framedtext}
\begin{choices}
  \item ㄱ, ㄴ
  \item ㄱ, ㄷ
  \item ㄴ, ㄷ
  \item ㄴ, ㄹ
  \item ㄷ, ㄹ
\end{choices}
\end{question}
\answer{2}
